\scanpage{47}\mspage{40}
% archive sheet number: 46
ne nous avance guère, \inserted{à lui \deleted{tout} seul} faute de savoir que l'axiome de factorisation se conserve, en passant de $(M,W\ldots)$ à $(M(I),W(I)\ldots)$.

C'est ici que le th. de factorisation du n° 1 \inserted{doit} s'avérer providentiel. Mais pour pouvoir l'appliquer aux couples $(C(I),C(I)_{*})$, $(TC(I),TC(I)_{*})$, il faut (en plus de l'accessibilité de $M$, ce qui est une condition fort modérée) des conditions de stabilité très fortes pour $C(I),TC(I)$.$^{*}$
% slanted
\margin{plus les \inserted{\uline{petit}s} syst. de générateurs (pour $C,TC$.}
À vrai dire, il \unsure{reste|peut} la stabilité de $C$ et de $C\cap W$ par limites inductives ordinales (ce qui implique la $=$ stabilité pour les $C(I)$, \deleted{$\unsure{C(I)\cap W}$} $(C\cap W)(I)=C(I)\cap W(I)$).

À dire vrai, je vois peu de chances de trouver une str. de $Qf$ $(C,TC,F,TF)$ associée à $W$, et satisfaisant \deleted{\unsure{aux}} les conditions dites y compris celles de stabilité de $C$ et de $C\cap W$ par $\varinjlim$ ordinales, sans supposer que $W$ elle-même ne soit déjà stable par \inserted{(petites)} limites inductives filtrantes.
